\documentclass[10pt]{article}

\usepackage{latexsym}
\usepackage{amssymb}
\usepackage{code}
\usepackage{graphicx}

\usepackage{listings}



\usepackage[english]{babel}
%\usepackage{blindtext}
%\usepackage{fontspec}
%\setmainfont{Arial}

%\renewcommand{\familydefault}{\sfdefault}
%\usepackage{blindtext}
\pagenumbering{arabic}  % Arabic page numbers GM July 2000

\usepackage{authoraftertitle}
\usepackage{fancyhdr}
% Clear the header and footer
\fancyhead{}
\fancyfoot{}
%\lhead{\includegraphics[height=0.7cm]{../logo/nyp-logo-int.png} }
\rhead{\scriptsize \MyTitle}
%\lfoot{\scriptsize Specialist Diploma in Business \& Big Data Analytics
%  \\ Copyright
 % \copyright\ Nanyang Polytechnic. All Rights Reserved.}
\rfoot{\thepage}
\pagestyle{fancy}


\renewcommand{\topfraction}{0.95}
\renewcommand{\textfraction}{0.02}
\renewcommand{\floatpagefraction}{0.95}


\bibliographystyle{plain}

%%\textwidth      150mm \textheight     210mm \oddsidemargin  -2mm \evensidemargin -2mm

\renewcommand{\baselinestretch}{0.987}

\setlength{\parskip}{0.0in}

%\input{macros-ms}

\newcommand{\Nturns}{\, \vdash_{\mbox{\scriptsize lnf}} \,}

\newcommand{\tr}[1]{}
%%\newcommand{\tr}[1]{#1}
%%\newcommand{\nottr}[1]{#1}
\newcommand{\nottr}[1]{}

%\newcommand{\implies}{\supset}
\newcommand{\clabel}[1]{\mbox{(#1)}}
\newcommand{\rat}[1]{\rightarrowtail_{#1}}
\newcommand{\arr}{\rightarrow}
\newcommand{\arrow}{\rightarrow}
\newcommand{\Arr}{\Rightarrow}
\newcommand{\atsign}{@}
\newcommand{\simparrow}[0]{\Longleftrightarrow}
\newcommand{\proparrow}[0]{\Longrightarrow}
\newcommand{\comment}[1]{}
\newcommand{\ignore}[1]{}
%\newcommand{\kl}[1]{{\bf KL:#1}} %%%{\marginpar{\sc kl}{\bf #1}}
\newcommand{\kl}[1]{}
\newcommand{\ms}[1]{{\bf MS:#1}}       %%{\marginpar{\sc ms}{\bf #1}}
%%\newcommand{\jw}[1]{\marginpar{\sc jw}{\bf #1}}
\newcommand{\pjs}[1]{}
%%\newcommand{\ms}[1]{}
\newcommand{\jw}[1]{}
%%\newcommand{\kl}[1]{}


\newcommand{\pow}{\^{}}
\newcommand{\venv}{\Delta}
\newcommand{\mleq}{\mbox{\tt leq}}
\newcommand{\mmleq}{\leq}
\newcommand{\mas}{\mbox{\tt as}}
\newcommand{\mfix}{\mu}

\newcommand{\mysection}[1]{\vspace*{-2mm}\section{#1}\vspace*{-1mm}}
\newcommand{\mysubsection}[1]{\vspace*{-1mm}\subsection{#1}\vspace*{-1mm}}

\newcounter{cnt}
\newtheorem{ex}{Example}
%\newenvironment{example}{
%        \begin{ex}\rm}%
%    {\hfill$\Box$\end{ex}}

\newenvironment{nexample}{
        \begin{ex}\rm}%
        {\end{ex}}

%%\newtheorem{rem}{Remark}
%%\renewenvironment{remark}{
%%        \begin{rem}\rm}%
%%    {\hfill$\Box$\end{rem}}

%%\newenvironment{myexample}{
%%        \begin{ex}\rm}%
%%  {\hfill $\Diamond$\end{ex}}
%\newenvironment{example}{
%        \begin{example}\rm}%
%    {\end{example}}

\newenvironment{ttline}{\begin{trivlist}\item \tt}{\end{trivlist}}
\newenvironment{ttprog}{\begin{trivlist}\small\item \tt
        \begin{tabbing}}{\end{tabbing}\end{trivlist}}


\newcommand{\figcode}[1]
        {\begin{figure*}[t]#1
        \end{figure*}}


\title{Introduction to Spark Machine Learning}
%\author{Kenny Zhuo Ming Lu\\
%  \multicolumn{1}{p{.7\textwidth}}{\centering
%  \emph{School of Information Technology \\ Nanyang Polytechnic \\
%    180 Ang Mo Kio Avenue 8, Singapore 569830}}
%}


\begin{document}
\maketitle \makeatactive
\thispagestyle{fancy}

%\lstset{language=scala}

%\begin{abstract}
%\end{abstract}



\section{Exercises }

\begin{enumerate}
   \item Start the Spark cluster
\begin{code}
$/opt/spark-1.5.2-bin-hadoop2-hive2-r/sbin/start-all.sh
\end{code}
%$
   \item We are going to transform some data from HDFS. 
   \item Start Hadoop. In the VM, start a terminal and run the following command
\begin{code}
$ start-dfs.sh && start-yarn.sh
\end{code} %$
   \item Check out the source codes.
           \begin{enumerate}
               \item If you find the {\tt \~{}/git/scala-bigdata-starter/} folder exists in your system, 
\begin{code}
$ cd ~/git/scala-bigdata-starter
$ git pull
\end{code}
               \item otherwise, clone it from github.
\begin{code}
$ cd ~
$ mkdir git
$ cd git 
$ git clone http://github.com/luzhuomi/scala-bigdata-starter.git
$ cd scala-bigdata-starter
\end{code}
%$
             \end{enumerate}
 \item Change to the following folder
\begin{code}
$ cd ~/git/scala-bigdata-starter/spark-ml-examples/
\end{code}
%$

\end{enumerate}

\subsection{Tweet classifier}

The school have collected some twitter data on the keyword ``jae''. 
However the term ``jae'' could refer to some 
Korean celebrities having ``jae'' as part of their names, or to joint
admission exercise in Singapore. 
The task here is to build a classifier to differentiate the KPOP tweets
mentioning ``jae'' as a person's name or otherwise.

For example, the following tweet message falls into the category of
Korean Pop because it seems talking about a person's name 
\begin{code}
crazy cool   jae s lee's pic of street singer reflected in raindrops
 tuesday on 2nd ave  
\end{code}
%
On the other hand, the following tweet is not revelant to KPOP. 
\begin{code}
accident closes jae valley rd drivers advised to avoid area seek 
alternate routes
\end{code}
%
To achieve the goal, we need to develop a classifier, which is a
supervised machine learning technique. In this example, we consider
using Support Vector Machine (SVM) as the classifier
algorithm. On the higher level, we need to ``train'' the model with
some manually labelled data and perform some tests against the trained
model. As part of the input requirement the SVM expect the input data
to represented as a label (either yes or no, 1 or 0) accompanied by
the feature vector. The feature vector is a vector of values which
uniquely differentiate one entry from another ideally.  In the machine
learning context, features have to be fixed by the programmers. 

For the ease of illustration, we use a vector of hash codes derived
from the words appearing in the tweet.  Consider the following source
code in {\tt src/main/scala/TweetSVMFilter.scala} 

\begin{code}
object TweetSVMFilter {

  val vector_fixed_size = 30 
  // fixed the size of each vector. 
  // if vectors have different sizes, the gradient descent algorithm will fail
  // cut off if it exceeds, pad zeros if it has less than 30 elements

  def hash(str:String):Int = str.toList.foldLeft(2147483647)((h,c) => 31*h + c)
  
  def to_words(tweet:String):List[String] = tweet.split(" ").toList

  def pad_cap(xs:List[Double],size:Int):List[Double] = xs match 
  {
    case Nil if size > 0 => List.fill(size)(0.0) // fill the rest with 0.0
    case Nil             => Nil 
    case (y::ys) if size == 0 => Nil
    case (y::ys)              => y::(pad_cap(ys,size-1))
  }

  def main(args: Array[String]) {
     
    // cluster mode
    val conf = new SparkConf().setAppName("Spark SVM")
    val sc = new SparkContext(conf)
    val posTXT:RDD[String] = sc
      .textFile("hdfs://127.0.0.1:9000/data/tweet/label_data/Kpop/*.txt") 
    val negTXT:RDD[String] = sc
      .textFile("hdfs://127.0.0.1:9000/data/tweet/label_data/othertweet/*.txt") 
    
    // convert the training data to labeled points
    val posLP:RDD[LabeledPoint] = posTXT.map( (twt:String) => 
    {
      val ws = to_words(twt).map(w => hash(w).toDouble)
      LabeledPoint(1.0, Vectors.dense(pad_cap(ws ,vector_fixed_size).toArray))
    })
    val negLP:RDD[LabeledPoint] = negTXT.map( (twt:String) => 
    { 
      val ws = to_words(twt).map(w => hash(w).toDouble)
      LabeledPoint(0.0, Vectors.dense(pad_cap(ws ,vector_fixed_size).toArray))
    })
    val data = negLP ++ posLP

    // Split data into training (60%) and test (40%).
    val splits = data.randomSplit(Array(0.6, 0.4), seed = 11L)
    val training = splits(0).cache()
    val test = splits(1)

    // Run training algorithm to build the model
    val numIterations = 100
    val model = SVMWithSGD.train(training, numIterations)

    // Clear the default threshold.
    model.clearThreshold()

    // Compute raw scores on the test set. 
    val scoreAndLabels = test.map { point =>
      val score = model.predict(point.features)
      (score, point.label)
    }

    // Get evaluation metrics.
    val metrics = new BinaryClassificationMetrics(scoreAndLabels)
    val auROC = metrics.areaUnderROC()

    println("Area under ROC = " + auROC)	
    sc.stop
  }
}
\end{code}
The {\tt hash} function hashes a string value into an integer.
The {\tt to\_words} function turns a sentence into words.
The {\tt pad\_cap} function takes a list of double values as input {\tt xs} and
a size {\tt size}. If the size of the input
list exceeds the given size, it truncates the input list down to the size.
If the size of the input falls below {\tt size}, it pads the remaining spaces with zeros.
This is to ensure all vectors to the SVM have the same size.
In the {\tt main} function, we first load the labeled data from the HDFS.


\end{document}
